\chapter{Sociedades, facciones y personajes}

\section{Facciones en el tablero de Sestao}

\subsection{Sociedad de Ciencias Naturales (SCN)}
Institución de larga tradición local, prestigiosa por sus estudios espeleológicos y paleontológicos. Aunque sus miembros son veteranos respetables, carecen de recursos económicos y viven eclipsados por el poder industrial. 
\begin{itemize}
  \item \lead{Francisco:} Paleontólogo obsesionado con los amonites biomecánicos. Se siente culpable por haber confiado en Kessler y busca redimirse ayudando a los chavales.
  \item \lead{Esteban Esteban:} Profesor de astronomía del Instituto Ángela Figuera. Entusiasta y bondadoso, organiza observaciones nocturnas y advierte sobre las anomalías del eclipse lunar del 8 de diciembre.
\end{itemize}

\kobfigure{03_sociedad_ciencias.jpg}{Francisco en su despacho de la Sociedad de Ciencias examinando los fósiles.}

\subsection{Babcock \& Wilcox y el Grupo Kessler}
Multinacional con peso institucional decisivo en la comarca. Dentro del complejo opera la división de investigación dirigida por el \lead{Dr. Victor Kessler}, secundado por ingenieros sin escrúpulos y un cuerpo de seguridad privada implacable. 
\begin{itemize}
  \item \lead{Dr. Victor Kessler:} Antagonista principal. Alto, esbelto, bata impecable y mirada gélida. Considera a Alex una mera batería biológica para conectar la red del Éter antes del eclipse.
  \item \lead{Guardias Jurados de B\&W:} Hombres uniformados con porras, perros y walkie-talkies en frecuencia interna que patrullan las naves y rastrean el pueblo buscando a Alex.
\end{itemize}

\subsection{Protección Civil y Policía Municipal}
Encargados del orden público y las emergencias cotidianas, a menudo desbordados por las protestas laborales de la reconversión.
\begin{itemize}
  \item \lead{Toño:} El voluntario más joven de Protección Civil. Conduce un viejo Seat Panda blanco y naranja; tiene buenas intenciones y siempre quiere ayudar, pero su torpeza puede meter a los PJs en líos.
\end{itemize}

\section{Tabla de rumores de Sestao (1d8)}

Tira 1d8 cuando los jugadores pregunten en los recreativos, en el colegio o a los parroquianos de las tabernas:

\begin{enumerate}[label=\textbf{\arabic*.}]
  \item \lead{La embarcación de La Benedicta:} Si bajas a la ría una noche de luna llena entre las 2 y las 3 de la mañana, ves una balsa de madera con adolescentes riendo que se tiran al agua: son los fantasmas del naufragio de hace un siglo. \emph{(Verdad a medias: son reflejos provocados por la energía del Éter en el agua)}.
  \item \lead{El laboratorio atómico de B\&W:} Dicen que en los sótanos de la vega de Galindo hay restos de uranio de la central de Lemóniz y que los científicos visten monos amarillos. \emph{(Cierto: investigan materiales prohibidos, pero es tecnología alienígena)}.
  \item \lead{La desaparición de Alex:} Hace tres años Alex se cayó a un pozo en La Arboleda y se lo tragó la tierra. Paco Lobatón sacó su foto en \emph{Quién sabe dónde}, pero la policía cerró el caso. \emph{(Totalmente cierto)}.
  \item \lead{Las luces de la Casa del Indiano:} En la buhardilla de las dos palmeras se reúnen miembros de una secta secreta para hacer cánticos raros. \emph{(Falso: son fugas del colector de Éter que resuena con los cimientos)}.
  \item \lead{Patxi el Bajo:} A un vecino bajito de la asociación de parados le ofrecieron un trabajo secreto en B\&W que le pagaría de por vida... y nunca volvió a aparecer. \emph{(Cierto: Kessler eliminó a los testigos que cavaron el túnel)}.
  \item \lead{El eclipse del 8 de diciembre:} El profesor Esteban dice que la luna se teñirá de rojo sangre y que las brújulas y relojes se volverán locos en toda Bizkaia. \emph{(Totalmente cierto)}.
  \item \lead{El monstruo de la mina:} Los mineros retirados de La Arboleda juran que en el fondo de las galerías hay conchas de piedra gigantes que laten con un corazón verde. \emph{(Cierto: la colonia de amonites cósmicos)}.
  \item \lead{Sirenas en la ría:} Unos pescadores de angulas aseguran haber visto una figura luminosa nadando a contracorriente sin respirar. \emph{(Cierto: Alex sintonizando sus poderes cerca del agua)}.
\end{enumerate}
